\chapter{12.11-3等边三角形综合}

\section*{一、选择题（共 9 小题）}

\questionlist[1]
\item
  关于等边三角形，下列说法中错误的是（\quad）
  \optionlist
  \item 等边三角形中，各边都相等
  \item 等边三角形是特殊的等腰三角形
  \item 三个角都等于 $60^{\circ}$ 的三角形是等边三角形
  \item 有一个角为 $60^{\circ}$ 的等腰三角形不是等边三角形
  \end{enumerate}

\item
  如图, $\triangle MNP$ 中, $\angle P=60^{\circ}$ , $MN=NP$ , $MQ \perp PN$ , 垂足为 $Q$ , 延长 $MN$ 至 $G$ , 取 $NG=NQ$ , 若 $\triangle MNP$ 的周长为 12, $MQ=a$ , 则 $\triangle MGQ$ 周长是 （\quad）
  \begin{questionwithimage}{2026-06-10-14-11_11-1.jpg}
  \optionlist
  \item $8+2a$
  \item $8+a$
  \item $6+a$
  \item $6+2a$
  \end{enumerate}
  \end{questionwithimage}

\item
  如图, 已知 $\angle ABC=120^{\circ}$ , $BD$ 平分 $\angle ABC$ , $\angle DAC=60^{\circ}$ , 若 $AB=2$ , $BC=3$ , 则 $BD$ 的长是 （\quad）
  \begin{questionwithimage}{2026-06-10-14-11_11-2.jpg}
  \optionlist
  \item 5
  \item 7
  \item 8
  \item 9
  \end{enumerate}
  \end{questionwithimage}

\item
  正三角形 $\triangle ABC$ 的边长为 3，依次在边 $AB$、$BC$、$CA$ 上取点 $A_{1}$、$B_{1}$、$C_{1}$，使 $AA_{1}=BB_{1}=CC_{1}=1$，则 $\triangle A_{1}B_{1}C_{1}$ 的面积是（\quad）
  \optionlist
  \item $\frac{\sqrt{3}}{4}$
  \item $\frac{3\sqrt{3}}{4}$
  \item $\frac{9}{4}$
  \item $\frac{9\sqrt{3}}{4}$
  \end{enumerate}

\item
  如图, $\triangle ABC$ 中, $\angle BAC=120^{\circ}$ , $AD \perp BC$ 于 $D$ , 且 $AB+BD=DC$ , 则 $\angle C$ 的大小是 （\quad）
  \begin{questionwithimage}{2026-06-10-14-11_11-3.jpg}
  \optionlist
  \item $20^{\circ}$
  \item $25^{\circ}$
  \item $30^{\circ}$
  \item 大于 $30^{\circ}$
  \end{enumerate}
  \end{questionwithimage}

\item
  如图,在 $Rt\triangle ABC$ , $\angle ACB=90^{\circ}$ , $CD$、 $CE$ 是斜边上的高和中线, $AC=CE=10cm$ , 则 $BD$ 长为 （\quad）
  \begin{questionwithimage}{2026-06-10-14-11_11-4.jpg}
  \optionlist
  \item $5cm$
  \item $10cm$
  \item $15cm$
  \item $25cm$ ,
  \end{enumerate}
  \end{questionwithimage}

\item
  如图, 六边形 $ABCDEF$ 中, 每一个内角都是 $120^{\circ}, AB=12, BC=30, CD=8, DE=28$ . 求这个六边形的周长为 （\quad）
  \begin{questionwithimage}{2026-06-10-14-11_11-5.jpg}
  \optionlist
  \item 125
  \item 126
  \item 116
  \item 108
  \end{enumerate}
  \end{questionwithimage}

\item
  如图, $P$ 是等边三角形 $ABC$ 内的一点, 且 $PA=3$ , $PB=4$ , $PC=5$ , 以 $BC$ 为边在 $\triangle ABC$ 外作 $\triangle BQC\cong \triangle BPA$ , 连接 $PQ$ , 则以下结论错误的是 （\quad）
  \begin{questionwithimage}{2026-06-10-14-11_11-6.jpg}
  \optionlist
  \item $\triangle BPQ$ 是等边三角形
  \item $\triangle PCQ$ 是直角三角形
  \item $\angle APB=150^\circ$
  \item $\angle APC=135^\circ$
  \end{enumerate}
  \end{questionwithimage}

\item
  如图所示, 在 $\triangle ABC$ 中, $AB=AC$ , $D$、 $E$ 是 $\triangle ABC$ 内两点, $AD$ 平分 $\angle BAC$ . $\angle EBC=\angle E=60^{\circ}$ , 若 $BE=6$ , $DE=2$ , 则 $BC$ 的长度是 （\quad）
  \begin{questionwithimage}{mineru-1.jpg}
  \optionlist
  \item 6
  \item 8
  \item 9
  \item 10
  \end{enumerate}
  \end{questionwithimage}

\end{enumerate}

\section*{二、填空题（共 5 小题）}

\blanklist[10]
\item
  如图, 在一个正方体的 2 个面上画了两条对角线 $AB$、 $AC$ , 那么这两条对角线的夹角是 \_\_\_\_.\underline{\hspace{2cm}}

  \centerquestionimage{0.35\textwidth}{mineru-2.jpg}

\item
  由于木质衣架没有柔性，在挂置衣服的时候不太方便操作。小敏设计了一种衣架，在使用时能轻易收拢，然后套进衣服后松开即可。如图\circled{1}，衣架杆 $OA=OB=18cm$，若衣架收拢时， $\angle AOB=60^{\circ}$，如图\circled{2}，则此时 $A, B$ 两点之间的距离是 \_\_\_\_ cm。\underline{\hspace{2cm}}

  \begin{questionimagegroup}
  \questionimageitem{0.32\textwidth}{mineru-3.jpg}{图\circled{1}}
  \questionimagegap
  \questionimageitem{0.32\textwidth}{mineru-4.jpg}{图\circled{2}}
  \end{questionimagegroup}

\item
  如图, $\triangle ABC$ 是等边三角形, 高 $AD$、$BE$ 相交于点 $H$ , $BC=4\sqrt{3}$ , 在 $BE$ 上截取 $BG=2$ , 以 $GE$ 为边作等边三角形 $GEF$ , 则 $\triangle ABH$ 与 $\triangle GEF$ 重叠 (阴影) 部分的面积为 \_\_\_\_.\underline{\hspace{2cm}}

  \centerquestionimage{0.4\textwidth}{mineru-5.jpg}

\item
  如图所示, 已知 $\angle 1=\angle 2$ , $AD=BD=4$ , $CE \perp AD$ , $2CE=AC$ , 那么 $DE$ 的长是 \_\_\_\_.\underline{\hspace{2cm}}

  \centerquestionimage{0.4\textwidth}{mineru-6.jpg}

\item
  如图, 直线 $a \parallel b$ , $\triangle ABC$ 是等边三角形, 点 $A$ 在直线 $a$ 上, 边 $BC$ 在直线 $b$ 上, 把 $\triangle ABC$ 沿 $BC$ 方向平移 $BC$ 的一半得到 $\triangle A'B'C'$ (如图\circled{1}); 继续以上的平移得到图\circled{2}, 再继续以上的平移得到图\circled{3}, ...; 请问在第 100 个图形中等边三角形的个数是 \_\_\_\_.\underline{\hspace{2cm}}

  \begin{questionimagegroup}
  \questionimageitem{0.28\textwidth}{mineru-7.jpg}{\circled{1}}
  \questionimagegap
  \questionimageitem{0.28\textwidth}{mineru-8.jpg}{\circled{2}}
  \questionimagegap
  \questionimageitem{0.28\textwidth}{mineru-9.jpg}{\circled{3}}
  \end{questionimagegroup}

\end{enumerate}

\section*{三、解答题（共 1 小题）}

\solutionlist[15]
\item
  如图, 已知 $\triangle ABC$ 为等边三角形, 延长 $BC$ 到 $D$ , 延长 $BA$ 到 $E$ , 并且使 $AE=BD$ , 连接 $CE$ , $DE$ . 求证: $EC=ED$ .

  \centerquestionimage{0.4\textwidth}{mineru-10.jpg}

\end{enumerate}
